


Generated by \texttt{python -m crashguard.r6\_bounds} and \texttt{python -m crashguard.r6\_docs\_bounds} from \texttt{results/results.json} (keys \texttt{r6.b.*}). Design: \texttt{docs/02\_preregistration/PREREGISTRATION\_V17.md} sections 7 and 8. Labels: ANALYTICAL for the formulas, TRACE\_DRIVEN for the stored trace results quoted next to them, MEASURED\_SOFTWARE for the sensing and L11 limits. None of this is a hypothesis test.

\subsection*{1. Exact bounds (R6-5)}

\textbf{Inputs.} Time on air by the Semtech formula as implemented in \texttt{crashguard.lora.airtime\_ms} (explicit header, CRC, coding rate 4/5, 125 kHz, 8 preamble symbols). Hourly allowance 1\% x 3,600 s = 36 s (the duty model of the ns-3 application and of the custom simulator; see section 3). The node learns that an attempt failed after the second receive window, 0.2 s after the end of the uplink, plus the 26-byte acknowledgement time. Retraction share rho = 0.1082 of provisional alarms (stored \texttt{ar.r3.a1.retraction\_share}). Frames per retracted event m\_r = 2.88 (the mean of 2.88 at SF7 and 2.89 at SF9, from the stored loss-free ns-3 outputs).

\textbf{Formulas.}
\begin{itemize}
\item Attempts that fit before a deadline D (back-to-back at the earliest feasible times): k(D) = floor((D - A97) / (A97 + 0.2 s + A26)) + 1.
\item Attempts per hour that the allowance pays for: floor(36 s / A).
\item Extra airtime of the two-stage alert, \textbf{lower bound} (no collisions, no retries; both only add airtime): LB = (A14 / A97) x (1 + rho (m\_r - 1) / (1 - rho)).
\item With the micro-confirm replacing the 97-byte confirm: (1 - rho)(A14 + A7) + rho (A14 + (m\_r - 1) A7), over (1 - rho) A97, minus one; and over R3's two-stage total.
\end{itemize}

{\scriptsize\raggedright
\begin{longtable}{p{0.086\linewidth}p{0.086\linewidth}p{0.086\linewidth}p{0.086\linewidth}p{0.086\linewidth}p{0.086\linewidth}p{0.086\linewidth}p{0.086\linewidth}p{0.086\linewidth}p{0.086\linewidth}p{0.086\linewidth}}
\toprule
SF & A7 & A14 & A26 & A97 (s) & attempts in 10 s & attempts in 30 s & A14/A97 & R3 bound & micro bound & micro over R3 \\
\midrule
\endhead
7 & 0.0361 & 0.0463 & 0.0617 & 0.1692 & 23 & 70 & 0.274 & 0.336 & -0.431 & 0.416 \\
8 & 0.0722 & 0.0824 & 0.1132 & 0.2975 & 16 & 49 & 0.277 & 0.340 & -0.391 & 0.443 \\
9 & 0.1239 & 0.1649 & 0.2058 & 0.5335 & 11 & 32 & 0.309 & 0.380 & -0.368 & 0.446 \\
10 & 0.2478 & 0.2888 & 0.4116 & 0.9851 & 6 & 19 & 0.293 & 0.360 & -0.362 & 0.457 \\
11 & 0.4956 & 0.6595 & 0.8233 & 2.1340 & 3 & 9 & 0.309 & 0.380 & -0.368 & 0.446 \\
12 & 0.9912 & 1.1551 & 1.6466 & 3.9404 & 2 & 5 & 0.293 & 0.360 & -0.362 & 0.457 \\
\bottomrule
\end{longtable}}

Columns "R3 bound" and "micro bound" are the extra airtime of the two-stage alert against the single-stage 97-byte alert; "micro over R3" is the airtime of the micro scheme per provisional alarm relative to R3's two-stage total.

\textbf{Reading.}
\begin{itemize}
\item The lower bound of R3's extra airtime is 0.336 at SF7 and 0.380 at SF9 and stays below 0.40 up to SF12 (\texttt{r6.b.rule.sf\_where\_lower\_bound\_reaches\_40pct} = 0, none). The measured SF9 value in ns-3 at loss 0 is 0.538: retries and collisions inflate the bound by a factor 1.42 at SF9 and 1.05 at SF7. The headroom to the 40\% bound (0.40 / bound) is 1.19 at SF7, 1.17 at SF8, 1.05 at SF9 and 1.11 at SF10: at SF9 only a 5\% inflation is tolerated, so the failure at SF9 is a failure of the retry and collision load, not of the frame sizes.
\item The micro-confirm's frame is 25\% shorter than the provisional frame at SF9 (A7 = 0.124 s against A14 = 0.165 s), and the micro scheme needs 0.45 of R3's airtime at SF9 (0.42 at SF7). Its extra airtime in the R3 definition is negative (-0.37 at SF9) because it drops the content of the 97-byte alert; the preregistration states this as a limit.
\item Tag length of the micro-confirm. A 5-byte tag has 2\textasciicircum{}-40 per attempt; at the attacker's ceiling of 997 frames per hour (SF7, 1\% allowance) the expected time to one accepted forgery is 1.26e+05 years if every frame hit a pending flash. A 4-byte tag would give 0.20\% per year and is rejected.
\end{itemize}

\textbf{Empirical boundaries (not bounds; intervals in the cited keys).}
\begin{itemize}
\item Sensing range: frontal-barrier pulses meet the 0.80 rule at +-32 g; pole pulses at +-64 g; no configuration meets it for side or other pulses (\texttt{p.b2.*}). Pole and side failed the generator's fidelity gate, so no claim about them is made. None of the estimators tested recovers pole and side impacts.
\item Trace finding: on links with loss of 40\% or more the no-hold control delivers 54.0\% and D1 49.3\% within 10 s; one alert uses at most 6.4 s, 17.8\% of the hourly allowance.
\item Mirror flips: accelerometer-only flips of x and y are detected in 57\% and 89\% of installations against 95\% required; both-channel flips in 100\%.
\end{itemize}

\subsection*{2. The spreading-factor rule (R6-7.2)}

The two-stage alert is enabled only where the headroom of the lower bound covers the retry and collision inflation. The inflation measured in ns-3 is 1.05 at SF7 (headroom 1.19: passes) and 1.42 at SF9 (headroom 1.05: fails), which agrees with the two cells that were run (SF7 pass, SF9 fail). SF8 has headroom 1.17 but was not run, and SF10 to SF12 have headroom 1.11, 1.05 and 1.11 and were not run. \textbf{Rule: enable the two-stage alert at SF7; at other spreading factors send the single-stage full decision, unless a run at that spreading factor shows an inflation below its headroom.} This is a design rule derived from the analysis and checked against two cells, not a tested hypothesis. The rule changes if the micro-confirm (R6-1) is used, because that scheme's extra airtime in the R3 definition is negative at every spreading factor (at the price of the dropped content). \textbf{Merged with the duty-lockout gate (delta verification run):} this rule is the load gate of one link-aware gating rule; the second gate, which matters only in a region that enforces a per-transmission off-time on one sub-band, is derived in \texttt{docs/03\_results/DELTA\_DUTY\_LOCKOUT.md} section 5. Under the hourly allowance evaluated here that gate is open at every spreading factor.

\subsection*{3. Duty model and the greedy-retry proposition (R6-7.1)}

\textbf{Which duty model the code uses.} The ns-3 application sets \texttt{g\_limit = duty x 3600} (xsim/ns3/xsim\_alert.cc) and a trailing-hour airtime window, and the custom simulator uses a trailing-window allowance (\texttt{world\_sim.Window}). It is an \textbf{hourly allowance}, not a per-transmission off-time T\_off = 99 x T\_on.

\textbf{Proposition (ANALYTICAL; dominant under the stated assumptions).} Let an alert be due at time 0 with deadline D. Assume (A1) the outcome of each attempt is independent of the time at which it is sent and has the same success probability q > 0; (A2) the earliest start of an attempt depends only on the end of the previous attempt, is non-decreasing in that end (the per-transmission off-time model T\_off = 99 x T\_on and the acknowledgement-wait rule both satisfy this) and an attempt must end by D; (A3) the allowance is not binding within the deadline. Then sending every attempt at its earliest feasible time maximises the number of attempts before D and therefore maximises the probability of delivery before D, so no spacing heuristic does better.

\emph{Proof.} Let f(t) be the earliest start of the next attempt when the previous attempt started at t; by (A2) f is non-decreasing. Let t\_1 < ... < t\_k be any feasible schedule and g\_1 = 0, g\_\{i+1\} = f(g\_i) the greedy schedule. Then g\_1 = 0 <= t\_1. If g\_i <= t\_i then g\_\{i+1\} = f(g\_i) <= f(t\_i) <= t\_\{i+1\} (the last step because t\_\{i+1\} is feasible after t\_i). So g\_i <= t\_i for every i, and the greedy attempt i also ends by D whenever attempt i of the other schedule does. Hence greedy fits at least as many attempts, say K\_g >= K. By (A1) the probability that at least one succeeds is 1 - (1 - q)\textasciicircum{}K, which is non-decreasing in K. (A3) holds in the stored runs: one alert uses at most 6.4 s, 17.8\% of the hourly allowance, and the exact bound above allows 23 attempts in 10 s at SF7 against 212 frames per hour that the allowance pays for. ? Every step above is an inequality between the quantities named; no step uses a property of the channel beyond (A1).

\textbf{Where it can fail.} Under bursty, time-correlated loss (A1 is false) a shifted schedule could land in a good state. The traces are the check: on the LoED held-out links with loss of 40\% or more, delivery within 10 s is 54.0\% for the no-hold control against 49.3\% for D1, and the control used at most 17.8\% of the hourly allowance. This is consistent with the proposition; it does not prove it, and the replay maps attempt j to the j-th next frame of the same device (a limit stated in the paper).

\subsection*{4. Quarantine indexing (R6-7.3)}

Read from the code. The only quarantine counter in the repository is \texttt{r5\_l8.process}; it is keyed on the vehicle and its state is (high-confidence forwards, last high slot, strikes, quarantined flag, last quarantine slot). The alert kinds it takes are \texttt{ok}, \texttt{implausible} and \texttt{repeat}; \textbf{there is no verification stage in the model}, so every alert attributed to a vehicle counts, including an over-rate alert, before any tag check. The implemented base verifier (\texttt{sct\_protocol.SctBase}, \texttt{packet.py}) checks the tag before anything else (\texttt{REJECT\_TAG}) and has no quarantine counter. Therefore: if the quarantine rule were implemented as the L8 model reads, keyed on the claimed vehicle id and counting before tag verification, an attacker without the victim's key could send frames with the victim's id (bad tags, repeated or implausible positions) and move the victim into quarantine. The bounded models of H-R6-3b decide it formally. \textbf{Design fix (two counters):} a per-cell counter for unverified frames, used only for flood defence and never quarantining a vehicle, and a per-vehicle rate limit applied only after successful verification.

